%2multibyte Version: 5.50.0.2960 CodePage: 1252

\documentclass{article}
%%%%%%%%%%%%%%%%%%%%%%%%%%%%%%%%%%%%%%%%%%%%%%%%%%%%%%%%%%%%%%%%%%%%%%%%%%%%%%%%%%%%%%%%%%%%%%%%%%%%%%%%%%%%%%%%%%%%%%%%%%%%%%%%%%%%%%%%%%%%%%%%%%%%%%%%%%%%%%%%%%%%%%%%%%%%%%%%%%%%%%%%%%%%%%%%%%%%%%%%%%%%%%%%%%%%%%%%%%%%%%%%%%%%%%%%%%%%%%%%%%%%%%%%%%%%
\usepackage{amsmath}

\setcounter{MaxMatrixCols}{10}
%TCIDATA{OutputFilter=LATEX.DLL}
%TCIDATA{Version=5.50.0.2960}
%TCIDATA{Codepage=1252}
%TCIDATA{<META NAME="SaveForMode" CONTENT="1">}
%TCIDATA{BibliographyScheme=Manual}
%TCIDATA{Created=Monday, September 22, 2003 09:40:11}
%TCIDATA{LastRevised=Thursday, September 25, 2025 13:21:40}
%TCIDATA{<META NAME="GraphicsSave" CONTENT="32">}
%TCIDATA{<META NAME="DocumentShell" CONTENT="General\SW\Blank - Standard LaTeX Article">}
%TCIDATA{Language=American English}
%TCIDATA{CSTFile=LaTeX article (bright).cst}

\newtheorem{theorem}{Theorem}
\newtheorem{acknowledgement}[theorem]{Acknowledgement}
\newtheorem{algorithm}[theorem]{Algorithm}
\newtheorem{axiom}[theorem]{Axiom}
\newtheorem{case}[theorem]{Case}
\newtheorem{claim}[theorem]{Claim}
\newtheorem{conclusion}[theorem]{Conclusion}
\newtheorem{condition}[theorem]{Condition}
\newtheorem{conjecture}[theorem]{Conjecture}
\newtheorem{corollary}[theorem]{Corollary}
\newtheorem{criterion}[theorem]{Criterion}
\newtheorem{definition}[theorem]{Definition}
\newtheorem{example}[theorem]{Example}
\newtheorem{exercise}[theorem]{Exercise}
\newtheorem{lemma}[theorem]{Lemma}
\newtheorem{notation}[theorem]{Notation}
\newtheorem{problem}[theorem]{Problem}
\newtheorem{proposition}[theorem]{Proposition}
\newtheorem{remark}[theorem]{Remark}
\newtheorem{solution}[theorem]{Solution}
\newtheorem{summary}[theorem]{Summary}
\newenvironment{proof}[1][Proof]{\textbf{#1.} }{\ \rule{0.5em}{0.5em}}
\input{tcilatex}
\input{tcilatex}
\setlength{\textwidth}{6.5in}
\setlength{\textheight}{9in}
\setlength{\topmargin}{-.50in}
\setlength{\oddsidemargin}{.0in}
\setlength{\evensidemargin}{.0in}

\begin{document}


\begin{center}
{\LARGE Assignment 3}
\end{center}

\noindent Blankenau\hspace*{\stretch{0.7500}}

\noindent Economics 905, Fall 2025

\noindent \textit{Instructions.} You may work in groups and may compare
answers prior to submission. Joint submissions are allowed and encouraged.
All work should be legible and concise. A due date will be announced in
class.

\begin{enumerate}
\item Consider a static economy with a single representative agent endowed
with one unit of time to allocate across labor and leisure and $k_{0}$ units
of capital. Preference are given by $U=\ln c+\beta \ln \ell $ where $\beta
>0,$ $c$ is the agent's consumption and $\ell $ is the agent's leisure$.$
There is also a representative firm with a production technology given by $%
y=k^{\alpha }n^{1-\alpha }\ $where $\alpha \in \left( 0,1\right) ,$ $k$ is
capital employed by the firm and $n$ is the labor employed by the firm$.$
Finally there is a government which consumes $G$ units of output and
finances its consumption by imposing a lump sum tax on the representative
agent subject to a balanced budget constraint. Government spending is a
fixed share of output (rather than a fixed quantity of output). That is, $%
G=gy$ so that the size of government depends on the size of output.

\begin{enumerate}
\item State the social planner's problem and solve the social planner's
problem for the optimal amount of $n$. The social planner does not choose $g$%
, but does choose $y$ and thus $G$.

\item Consider a competitive equilibrium.

\begin{enumerate}
\item State the agents's problem and find the first order conditions for the
agents's problem.

\item State the firm's problem and find the first order conditions for the
firm's problem.

\item Find the government budget constraint.

\item Find $n$ in a competitive equilibrium.
\end{enumerate}

\item Now suppose the situation is the same, except the representative agent
values the government spending such that preferences are $U=\ln c+\beta \ln
\ell +\theta \ln G$ where still $G=gy$ and $G$ is still taken as given by
the agent. In this setting the social planner chooses $g$. State the social
planner's problem and solve the social planner's problem for the optimal
amount of $n$ and $g$.
\end{enumerate}

\item Consider an economy with a continuum of consumers and normalize the
mass of consumers to 1. Each agent is endowed with one unit of time which is
allocated across work and leisure. Preferences are $U\left( c,\ell ,\bar{\ell%
}\right) =\ln c+\ln \ell +\ln \bar{\ell}$ where $c$ and $\ell $ are the
individual's consumption and leisure and $\bar{\ell}$ is the average leisure
across the population. Note that because any individual is very small
relative to the population, each consumer will treat $\bar{\ell}$ as given.
There is an externality in that the utility of an individual is influenced
by others' leisure. A representative firm behaves competitively and hires
labor, $n$, to maximize profits with production technology given by $%
y=n^{\alpha }$.

\begin{enumerate}
\item Determine the Pareto optimal leisure.

\item Determine leisure in a competitive equilibrium.

\item Now suppose that the government taxes labor income at rate $\tau $ and
returns the revenue lump sum to the representative agents as a subsidy, $%
\varsigma $, subject to a balanced budget constraint. Write down the agent's
optimization problem and find the first order condition (or conditions
depending on how you have written the problem). Find and impose any
equilibrium conditions that are required so that you can solve for $\ell $
in a competitive equilibrium. Solve for $\ell .$

\item Find the level of taxation that leads to the Pareto optimal allocation.
\end{enumerate}

\item Consider an economy with many agents, each of whom maximizes 
\begin{equation*}
\ln c
\end{equation*}%
where $c$ is consumption. The representative agent has one unit of time
available in each period and must accumulate human capital before he works.
After getting human capital, the agent produces according to $c=\alpha hu$
where $\alpha >0,$ $h$ is the agent's human capital and $u$ is time spent in
consumption goods production by the agent, so that $s=\left( 1-u\right) $
units of time are spent accumulating human capital. The agent accumulates
human capital according to $h=zs^{\theta }\bar{s}^{1-\theta }$ where $%
0<\theta <1.$ Here $\bar{s}$ is the average amount of time spent studying
and is taken as given by the agent.

\begin{enumerate}
\item Find the solution to the social planner's problem. In particular find
the optimal level of $u$.

\item Find the allocation of time that would arise in a competitive
equilibrium.

\item Now suppose that the government pays people to go to school and taxes
them to pay for the subsidies. Specifically the government pays $\varsigma w$
per units of time in school where $w$ is the \textquotedblleft
wage\textquotedblright\ per unit of human capital (what a firm would pay to
hire one unit of time; i.e. $w=\alpha h$) and $\varsigma <1$ is the rate at
which schooling is subsidized$.$ The subsidy is paid for by a lump sum tax
and government must balance its budget. Find the allocation of time that
would arise in a competitive equilibrium. Find the optimal subsidy.

\item For the remainder of this problem there is no tax or subsidy. The
environment is as above except utility is given by $\ln c+\beta \ln \ell $
where $\ell $ is leisure. The representative agent still has one unit of
time available in each period and must accumulate human capital before
working. Thus $u+h+\ell =1.$ Again $h=zs^{\theta }\bar{s}^{1-\theta }$ where 
$s$ and $\bar{s}$ are again time and average time spent in education but now 
$s\neq (1-u)$ since some time is spent on leisure. After getting human
capital, the agent produces according to $c=\alpha hu$ as before. Find the
level of $u$ and $\ell $\ that arise as the solution to the social planner's
problem. Hint: Another endogenous item will require another first order
condition. \pagebreak
\end{enumerate}
\end{enumerate}

\begin{enumerate}
\item ..

\begin{enumerate}
\item 
\begin{equation*}
L=\ln (\left( 1-g\right) k^{\alpha }n^{1-\alpha })+\beta \ln \left(
1-n\right)
\end{equation*}%
First order conditions give 
\begin{equation*}
\frac{\left( 1-g\right) \left( 1-\alpha \right) k^{\alpha }n^{-\alpha }}{%
\left( 1-g\right) k^{\alpha }n^{1-\alpha }}=\frac{\beta }{1-n}
\end{equation*}%
solving gives%
\begin{eqnarray*}
\frac{1-\alpha }{n} &=&\frac{\beta }{1-n} \\
n &=&\frac{1-\alpha }{1-\alpha +\beta }
\end{eqnarray*}

\item $\allowbreak $

\begin{enumerate}
\item Agent problem is to max%
\begin{equation*}
\ln \left( wn+rk_{0}-\tau \right) +\beta \ln \left( 1-n\right)
\end{equation*}%
First order conditions give%
\begin{equation*}
\frac{w}{wn+rk_{0}-\tau }=\frac{\beta }{1-n}
\end{equation*}

\item Write down firms' optimization problem. On your own. The firm's
maximization gives 
\begin{equation*}
r=\alpha \left( \frac{n}{k}\right) ^{1-\alpha }
\end{equation*}%
\begin{equation*}
w=\left( 1-\alpha \right) \left( \frac{k}{n}\right) ^{\alpha }
\end{equation*}

\item Government budget constraint gives $G=\tau $ or $gy=\tau .$

\item Use market clearing conditions to find\ 
\begin{equation*}
\frac{\left( 1-\alpha \right) \left( \frac{k}{n}\right) ^{\alpha }}{%
k^{\alpha }n^{1-\alpha }-G}=\frac{\beta }{1-n}
\end{equation*}%
but $G=gy$ so we get 
\begin{equation*}
\frac{\left( 1-\alpha \right) k_{0}^{\alpha }n^{-\alpha }}{\left( 1-g\right)
k_{0}^{\alpha }n^{1-\alpha }}=\frac{\beta }{1-n}
\end{equation*}%
and simplifying gives 
\begin{equation*}
\frac{\left( 1-\alpha \right) }{\left( 1-g\right) n}=\frac{\beta }{1-n}
\end{equation*}%
or 
\begin{equation*}
\frac{\left( 1-\alpha \right) }{1+\beta \left( 1-g\right) -\alpha }=n
\end{equation*}
\end{enumerate}

\item ..

\begin{enumerate}
\item problem is to max%
\begin{equation*}
\ln (\left( 1-g\right) k^{\alpha }n^{1-\alpha })+\beta \ln \left( 1-n\right)
+\theta \ln gk^{\alpha }n^{1-\alpha }
\end{equation*}%
First order conditions give 
\begin{eqnarray*}
\frac{\left( 1-g\right) \left( 1-\alpha \right) k^{\alpha }n^{-\alpha }}{%
\left( 1-g\right) k^{\alpha }n^{1-\alpha }}+\theta \frac{\left( 1-\alpha
\right) gk^{\alpha }n^{-\alpha }}{gk^{\alpha }n^{1-\alpha }} &=&\frac{\beta 
}{1-n} \\
\frac{1}{1-g} &=&\frac{\theta }{g}
\end{eqnarray*}%
The first gives%
\begin{eqnarray*}
\frac{\left( 1-\alpha \right) }{n}+\theta \frac{\left( 1-\alpha \right) }{n}
&=&\frac{\beta }{1-n} \\
\left( 1-n\right) \left( \left( 1-\alpha \right) +\theta \left( 1-\alpha
\right) \right) &=&\beta n
\end{eqnarray*}%
\begin{equation*}
\frac{\left( 1-\alpha \right) +\theta \left( 1-\alpha \right) }{\beta
+\left( 1-\alpha \right) +\theta \left( 1-\alpha \right) }=n
\end{equation*}%
solve second gives$.$ 
\begin{equation*}
g=\frac{\theta }{1+\theta }
\end{equation*}
\end{enumerate}
\end{enumerate}

\item ...

\begin{enumerate}
\item Social planner's problem Max:%
\begin{equation*}
U\left( c,\ell ,\bar{c}\right) =\ln c+\ln \ell +\ln \bar{\ell}
\end{equation*}%
subject to $c=\left( 1-\ell \right) ^{\alpha }.$ Internalizes externality so 
\begin{equation*}
U\left( c,\ell ,\bar{\ell}\right) =\ln \left( 1-\ell \right) ^{\alpha }+2\ln
\ell 
\end{equation*}%
\begin{equation*}
\frac{\alpha }{1-\ell }=\frac{2}{\ell }
\end{equation*}%
$\allowbreak $%
\begin{equation*}
\ell =\frac{2}{\alpha +2}
\end{equation*}

\item Now look at competitive equilibrium. \ Agent's problem is 
\begin{equation*}
U\left( c,\ell ,\bar{c}\right) =\ln c+\ln \ell +\ln \bar{\ell}
\end{equation*}%
st. $c=\omega n$ 
\begin{equation*}
U\left( c,\ell ,\bar{\ell}\right) =\ln \left( w\left( 1-\ell \right) \right)
+\ln \ell +\ln \bar{\ell}
\end{equation*}%
\begin{equation*}
\frac{1}{1-\ell }=\frac{1}{\ell }
\end{equation*}%
\begin{equation*}
\ell =\frac{1}{2}
\end{equation*}

\item Under the subsidy scheme, the budget constraint in%
\begin{equation*}
c=\omega n\left( 1-\tau \right) +\varsigma 
\end{equation*}%
\begin{equation*}
U\left( c,\ell ,\bar{\ell}\right) =\ln \left[ \omega \left( 1-\ell \right)
\left( 1-\tau \right) +\varsigma \right] +\ln \ell +\ln \bar{\ell}
\end{equation*}%
\begin{equation*}
\frac{\omega \left( 1-\tau \right) }{\omega \left( 1-\ell \right) \left(
1-\tau \right) +\varsigma }=\frac{1}{\ell }
\end{equation*}%
Government budget constraint is $\omega \left( 1-\ell \right) \tau
=\varsigma .$ Wage is $\omega =\alpha \left( 1-\ell \right) ^{\alpha -1}\ $%
giving%
\begin{equation*}
\frac{\omega \left( 1-\tau \right) }{\omega \left( 1-\ell \right) \left(
1-\tau \right) +\omega \left( 1-\ell \right) \tau }=\frac{1}{\ell }
\end{equation*}%
\begin{equation*}
\frac{1-\tau }{1-\ell }=\frac{1}{\ell }
\end{equation*}%
\begin{equation*}
\ell =\frac{1}{2-\tau }
\end{equation*}

\item Set $\tau $ to make these equal. 
\begin{equation*}
\frac{2}{\alpha +2}=\frac{1}{1-\tau +1}
\end{equation*}%
\begin{equation*}
\frac{2}{\alpha +2}=\frac{1}{2-\tau }
\end{equation*}%
\begin{equation*}
\tau =1-\frac{1}{2}\alpha 
\end{equation*}
\end{enumerate}

\item ...

\begin{enumerate}
\item The social planner's problem is 
\begin{equation*}
\underset{c,u,h}{\max }\ln c
\end{equation*}%
subject to 
\begin{eqnarray*}
c &=&\alpha hu \\
h &=&z(1-u)
\end{eqnarray*}%
Rewrite as%
\begin{equation*}
\underset{u}{\max }\ln \alpha z(1-u)u
\end{equation*}%
This is the same a maximizing%
\begin{equation*}
(1-u)u
\end{equation*}%
The FOC is 
\begin{equation*}
1=2u
\end{equation*}%
so%
\begin{equation*}
u^{\ast }=\frac{1}{2}
\end{equation*}

\item An agents maximizes%
\begin{equation*}
\underset{c,u,h}{\ln }x
\end{equation*}%
subject to 
\begin{equation*}
\begin{array}{ll}
c=\alpha hu & h=z(1-u)^{\theta }\left( 1-\bar{u}\right) ^{\left( 1-\theta
\right) }%
\end{array}%
\end{equation*}%
Rewrite as%
\begin{equation*}
\underset{u}{\max }\ln \alpha z(1-u)^{\theta }\left( 1-\bar{u}\right)
^{\left( 1-\theta \right) }u
\end{equation*}%
Since the agent takes $\overline{u}$ as given this is the same as maximizing%
\begin{equation*}
(1-u)^{\theta }u
\end{equation*}%
The first order condition is%
\begin{equation*}
\theta (1-u)^{\theta -1}u=(1-u)^{\theta }
\end{equation*}%
\begin{equation*}
\theta u=1-u
\end{equation*}%
giving%
\begin{equation*}
u=\frac{1}{1+\theta }
\end{equation*}

\item Now the agent's problem is 
\begin{equation*}
\underset{u}{\max }\ln \left( \alpha z(1-u)^{\theta }\left( 1-\bar{u}\right)
^{\left( 1-\theta \right) }u\right) +(1-u)\varsigma w-\tau
\end{equation*}%
which is maximized at 
\begin{equation*}
\theta \alpha z(1-u)^{\theta -1}\left( 1-\bar{u}\right) ^{\left( 1-\theta
\right) }u+\theta \alpha z(1-u)^{\theta -1}\left( 1-\bar{u}\right) ^{\left(
1-\theta \right) }=\varsigma w
\end{equation*}%
In an equilibrium $w=\alpha z(1-u)^{\theta }\left( 1-\bar{u}\right) ^{\left(
1-\theta \right) }$ and $\overline{u}=u.$ This becomes%
\begin{equation*}
-\theta \alpha z(1-u)^{\theta -1}\left( 1-\bar{u}\right) ^{\left( 1-\theta
\right) }u+\alpha z(1-u)^{\theta }\left( 1-\bar{u}\right) ^{\left( 1-\theta
\right) }=\varsigma \alpha z(1-u)^{\theta }\left( 1-\bar{u}\right) ^{\left(
1-\theta \right) }
\end{equation*}%
\begin{equation*}
-\theta \alpha zu+\alpha z(1-u)=(1-u)\varsigma \alpha z
\end{equation*}%
\begin{equation*}
-\theta u+(1-u)=(1-u)\varsigma
\end{equation*}%
$\allowbreak $The solution is 
\begin{equation*}
u=\frac{1-\varsigma }{\theta +1-\varsigma }
\end{equation*}%
The optimal subsidy will be such that 
\begin{equation*}
\frac{1-\varsigma }{\theta +1-\varsigma }=\frac{1}{2}
\end{equation*}%
giving%
\begin{equation*}
\varsigma =1-\theta
\end{equation*}%
\begin{equation*}
\frac{1-\left( 1-\theta \right) }{\theta +1-\left( 1-\theta \right) }-\frac{1%
}{2}
\end{equation*}

\item The social planner's problem is 
\begin{equation*}
\underset{c,u,h,\ell }{\max }\ln c+\beta \ln \ell
\end{equation*}%
subject to 
\begin{eqnarray*}
c &=&\alpha hu \\
h &=&z(1-u-\ell )
\end{eqnarray*}%
Rewrite as%
\begin{equation*}
\underset{u}{\max }\ln \alpha z(1-u-\ell )u+\beta \ln \ell )
\end{equation*}%
The FOCs are 
\begin{equation*}
\frac{1-2u-\ell }{\left( 1-u-\ell \right) u}=0
\end{equation*}%
\begin{equation*}
\frac{1}{1-u-\ell }=\frac{\beta }{\ell }
\end{equation*}%
From the first 
\begin{equation*}
\ell =1-2u
\end{equation*}%
This into the second gives 
\begin{equation*}
\frac{1}{1-u-1-2u}=\frac{\beta }{1-2u}
\end{equation*}%
\begin{equation*}
\frac{1}{1-u-\left( 1-2u\right) }=\frac{\beta }{1-2u}
\end{equation*}%
\begin{equation*}
u=\frac{1}{\beta +2}
\end{equation*}%
\begin{equation*}
\ell =1-2\frac{1}{\beta +2}=\frac{\beta }{\beta +2}
\end{equation*}
\end{enumerate}
\end{enumerate}

\end{document}
